\subsection{Nakayama's Lemma}

PROPOSITION 6. --- For every A-module M, we have \(\Re(A)M \subset \Re(M)\) . We have equality if the A-module M is projective.

Let P be a maximal submodule of M; the A-module M/P is simple, hence annihilated by \(\Re(\mathrm{A})\) by Proposition 5 of VIII, p.~155. We therefore have \(\Re(\mathrm{A})\mathrm{M}\subset\mathrm{P}\) for every maximal submodule P of M, hence \(\Re(\mathrm{A})\mathrm{M}\subset\Re(\mathrm{M})\) .

We clearly have \(\mathfrak{R}(A_{s})=\mathfrak{R}(A)A_{s}\) . If the A-module M is projective, then there exists an A-module N such that \(M\oplus N\) is free, that is, the direct sum of a family \((\mathrm{L}_{i})_{i\in\mathrm{I}}\) of modules isomorphic to \(A_{s}\) . By Corollary 2 of VIII, p.~152, we have \(\mathfrak{R}(M\oplus N)=\mathfrak{R}(M)\oplus\mathfrak{R}(N)\) and \(\mathfrak{R}\big(\bigoplus_{i\in I}L_{i}\big)=\bigoplus_{i\in I}\mathfrak{R}(L_{i})\) . We then deduce \(\mathfrak{R}(M)=\mathfrak{R}(A)M\) from the equality \(\mathfrak{R}(L_{i})=\mathfrak{R}(A)L_{i}\) .

THEOREM 2 (``Nakayama's lemma''). --- Let M be an A-module and a two-sided ideal of A. Suppose that one of the following conditions is satisfied:

\begin{enumerate}
\def\labelenumi{(\roman{enumi})}
\item
  The A-module M is finitely generated, and \(\mathfrak{a}\) is contained in the radical of A.
\item
  The ideal \(\mathfrak{a}\) is nilpotent.
\end{enumerate}

If N is a submodule of M such that \(M = N + aM\) , then we have N = M. In particular, if the module M is nonzero, then we have \(M \neq aM\) .

Suppose that M is finitely generated and that we have \(\mathfrak{a} \subset \Re(\mathrm{A})\) . Let N be a submodule of M such that \(M = N + aM\) . By Proposition 6, we have \(M = N + \Re(M)\) , hence M = N by Proposition 2, b) of VIII, p.~153.

Now suppose that a is nilpotent, and let N be a submodule of M such that \(M = N + aM\) . By induction on the integer \(n \geqslant 0\) , we establish the relation \(M = N + a^{n}M\) . By assumption, there exists an integer \(n \geqslant 0\) such that \(a^{n} = 0\) ; hence M = N.

The last assertion of the theorem follows from the above by taking N to be 0.

COROLLARY 1. --- We keep the assumptions of Theorem 2. Let \((x_{i})_{i\in\mathrm{I}}\) be a family of elements of M, and let \(\overline{x}_{i}\) be the canonical image of \(x_{i}\) in M/aM. If the family \((\overline{x}_{i})_{i\in\mathrm{I}}\) generates the (A/a)-module M/aM, then the family \((x_{i})_{i\in\mathrm{I}}\) generates the A-module M.

This follows from Theorem 2 applied to the submodule N of M generated by the family \((x_{i})_{i\in \mathbf{I}}\) .

COROLLARY 2. --- We keep the assumptions of Theorem 2. Furthermore, let M' be an A-module and u : M' → M be a homomorphism. If the homomorphism \(\overline{u}\) from M'/aM' to M/aM deduced from u by passing to the quotients is surjective, then u is surjective.

It suffices to apply Theorem 2 to the image N of u: indeed, the image of \(\overline{u}\) is \((\mathrm{N} + \mathfrak{aM})/\mathfrak{aM}\) , so \(\overline{u}\) is surjective if and only if we have \(N + aM = M\) .

